\chapter{Conclus\~{a}o e Trabalho Futuro} \label{cap:conclusao}

\section{Conclus\~{a}o} \label{sec:conclusao}

O projeto \textit{IPW} teve como objetivo principal conceber e desenvolver uma plataforma \textit{web} para apoiar a gest\~{a}o de processos de averigua\c{c}\~{a}o no contexto da CAPT Consulting. A solu\c{c}\~{a}o proposta procura responder \`{a}s limita\c{c}\~{o}es identificadas no processo atual, nomeadamente a dispers\~{a}o da informa\c{c}\~{a}o, a dificuldade em acompanhar o estado dos processos e a inexist\^{e}ncia de um \textit{workflow} centralizado e estruturado.

Ao longo do desenvolvimento foi definida uma arquitetura composta por uma aplica\c{c}\~{a}o \textit{frontend}, uma \textit{API} \textit{backend} e uma base de dados relacional. O \textit{frontend} disponibiliza a interface de utiliza\c{c}\~{a}o da plataforma, enquanto o \textit{backend} concentra a l\'{o}gica de neg\'{o}cio, a autentica\c{c}\~{a}o, a autoriza\c{c}\~{a}o e o acesso aos dados. A base de dados foi modelada de forma a representar os principais conceitos do dom\'{i}nio, incluindo utilizadores, pap\'{e}is, \'{a}reas, processos, estados e hist\'{o}rico.

Na fase atual, a aplica\c{c}\~{a}o j\'{a} inclui funcionalidades fundamentais, como autentica\c{c}\~{a}o de utilizadores, sele\c{c}\~{a}o de papel, prote\c{c}\~{a}o de rotas, gest\~{a}o de utilizadores, cria\c{c}\~{a}o e listagem de processos, associa\c{c}\~{a}o de intervenientes e consulta de informa\c{c}\~{a}o relevante para o acompanhamento dos processos. Estas funcionalidades permitem validar a estrutura base da solu\c{c}\~{a}o e demonstrar a viabilidade da abordagem adotada.

O desenvolvimento permitiu tamb\'{e}m consolidar decis\~{o}es t\'{e}cnicas importantes, como a utiliza\c{c}\~{a}o de \textit{JWT} para autentica\c{c}\~{a}o, a aplica\c{c}\~{a}o de controlo de acessos baseado em pap\'{e}is e a separa\c{c}\~{a}o do \textit{backend} em camadas com responsabilidades distintas. Apesar de ainda existirem funcionalidades por concluir, a vers\~{a}o atual estabelece uma base consistente para a evolu\c{c}\~{a}o do sistema.

\section{Trabalho futuro} \label{sec:trabalho-futuro}

Como trabalho futuro, uma das funcionalidades mais relevantes a desenvolver \'{e} o suporte \`{a} anexa\c{c}\~{a}o de ficheiros aos processos. Esta funcionalidade permitir\'{a} associar fotografias, v\'{i}deos e outros documentos relevantes ao processo de averigua\c{c}\~{a}o. Para este efeito, dever\'{a} ser avaliada uma solu\c{c}\~{a}o de armazenamento de objetos, mantendo na base de dados apenas os metadados dos ficheiros e uma refer\^{e}ncia para a sua localiza\c{c}\~{a}o.

Outra melhoria prevista consiste no aprofundamento da integra\c{c}\~{a}o entre o \textit{frontend} e o \textit{backend}. Embora a aplica\c{c}\~{a}o j\'{a} possua v\'{a}rios componentes e \textit{endpoints} implementados, ainda existem \'{a}reas que necessitam de maior integra\c{c}\~{a}o, valida\c{c}\~{a}o e tratamento de estados de erro. Esta evolu\c{c}\~{a}o permitir\'{a} aproximar a experi\^{e}ncia de utiliza\c{c}\~{a}o do fluxo final pretendido.

Est\'{a} tamb\'{e}m prevista a melhoria dos mecanismos de monitoriza\c{c}\~{a}o e visualiza\c{c}\~{a}o de indicadores operacionais. Esta evolu\c{c}\~{a}o poder\'{a} incluir informa\c{c}\~{a}o agregada sobre prazos, distribui\c{c}\~{a}o de trabalho por utilizador ou \'{a}rea, estados dos processos e outros indicadores relevantes para a gest\~{a}o.

Adicionalmente, poder\'{a} ser estudada a atribui\c{c}\~{a}o autom\'{a}tica de processos a averiguadores e supervisores, tendo em conta crit\'{e}rios como \'{a}rea de especializa\c{c}\~{a}o, disponibilidade ou carga de trabalho. Esta funcionalidade foi considerada como requisito opcional, uma vez que depende da estabiliza\c{c}\~{a}o pr\'{e}via do fluxo principal de cria\c{c}\~{a}o e acompanhamento de processos.

Por fim, a integra\c{c}\~{a}o com o sistema IRA constitui uma possibilidade de evolu\c{c}\~{a}o futura. Esta integra\c{c}\~{a}o permitiria comunicar com outro sistema desenvolvido para o mesmo cliente, mas exige a defini\c{c}\~{a}o clara dos contratos de comunica\c{c}\~{a}o entre plataformas. Assim, esta funcionalidade dever\'{a} ser considerada ap\'{o}s a consolida\c{c}\~{a}o do n\'{u}cleo principal da aplica\c{c}\~{a}o.
